\section{Certified phase balls and their atlas}
\label{sec:phase-cover}

This section converts a cell retained by the completion search of
Section~\ref{sec:completion} into a ball of real phase matrices.  For a
center \(P\), the matrix \(e(P)\) need not be an exact Hadamard matrix.
The certificates control the center's residual and the displacement of
every exact solution in the ball.  Phase matrices and their ball radii
are measured in cycles.  Angular displacements explicitly identified
below are measured in radians.

We use the entrywise exponential \(e\) and the orthogonal double-centering
projection \(\Pi\) defined in the seed-cover section.  Real \(6\)-by-\(6\)
matrices are identified with \(\R^{36}\) in row-major order.  A phase
displacement is \emph{horizontal} if it belongs to \(\operatorname{ran}\Pi\).
All vector norms are Euclidean, all phase-matrix norms are Frobenius, and
norms of linear maps are induced Euclidean norms unless a subscript
specifies otherwise.  Stars on real matrices mean transpose.

\subsection{Residual and Root certificates}
\label{sec:root-certificate}

For a real phase matrix \(P\), define \(F(P)\in\C^{15}\) by
\begin{equation}
 F_{ij}(P)=\frac{1}{2\pi i}\sum_{g=0}^{5}e(P_{ig}-P_{jg}),
 \qquad 0\leq j<i\leq5.
 \label{phase:residual}
\end{equation}
Edges are ordered by increasing \(i\), then increasing \(j\); a complex
coordinate is stored as its real part followed by its imaginary part.
We also regard \(F(P)\) as a real \(30\)-vector.  Since every entry of
\(e(P)\) has modulus one, \(F(P)=0\) is equivalent to the row Hadamard
conditions, and hence to all the Hadamard conditions.  Its real
derivative \(L_P:\R^{36}\longrightarrow\R^{30}\), in complex notation, is
\begin{equation}
 (L_Px)_{ij}=\sum_{g=0}^{5}e(P_{ig}-P_{jg})(x_{ig}-x_{jg}).
 \label{phase:derivative}
\end{equation}
The factor \(2\pi i\) from differentiation cancels the denominator in
\eqref{phase:residual}.  These are the residual and full derivative in
\texttt{buildH::gen}, in \texttt{patchx.cpp}.

\begin{definition}[Root certificate]
\label{phase:root-data}
A \texttt{Root} consists of a stored real center \(P\), a stored real
\(36\)-by-\(k\) matrix \(T\), and nonnegative constants
\[
 \mathtt{C},\quad \mathtt{n0},\quad \mathtt{id},\quad
 \mathtt{ort},\quad \mathtt{hor},\quad
 \mathtt{mnorm},\quad \mathtt{froot}.
\]
The source stores column \(a\) of \(T\) as \texttt{T[a]} and uses
\(0\leq k\leq10\).  Write
\begin{equation}
 n_P(s)=\mathtt{n0}+\mathtt{id}\,s+\mathtt{C}\,s^2,
 \qquad s\geq0.
 \label{phase:normal-polynomial}
\end{equation}
It is certified if the basic bounds
\begin{equation}
 \|T^*T-I\|\leq\mathtt{ort},\qquad
 \|(I-\Pi)T\|\leq\mathtt{hor},\qquad
 \|F(P)\|\leq\mathtt{froot}
 \label{phase:root-basic}
\end{equation}
hold, together with the following displacement assertions.  For every
horizontal \(\delta\) satisfying \(F(P+\delta)=0\), put
\(\eta=T^*\delta\) and \(s=\|\delta\|\).  Then
\begin{align}
 \|(I-TT^*)\delta\|&\leq n_P(s), \label{phase:root-normal}\\
 s^2&\leq(1+\mathtt{ort})\|\eta\|^2+
                  \|(I-TT^*)\delta\|^2. \label{phase:root-square}
\end{align}
For every such \(\delta\) and every horizontal \(v\), with
\(h'=\delta-v\), one also has
\begin{equation}
 \|(I-TT^*)h'\|
 \leq\bigl[\mathtt{id}+\mathtt C(\|\delta\|+\|v\|)\bigr]\|h'\|
       +\mathtt{mnorm}\,\|F(P+v)\|.
 \label{phase:root-difference}
\end{equation}
These assertions have no radius restriction.  A Root certificate
neither requires \(F(P)=0\) nor asserts the existence of an exact
solution.  A radius will be attached when a phase ball is recorded.
\end{definition}

The first two displacement bounds explain how the certificate will be
used.  For such a \(\delta\), if \(s_0\geq\|\delta\|\) and
\(t_0\geq\|T^*\delta\|\), then nonnegativity of the coefficients of
\(n_P\) gives
\[
 \|\delta\|\leq
 \left(n_P(s_0)^2+(1+\mathtt{ort})t_0^2\right)^{1/2}.
\]
For a retained completion cell we shall bound both the full displacement
and its tangent coordinates.  The displayed bound can then improve the
radius that covers all exact completions in the cell.  The difference
bound will compare displacements from two centers during refinement.

\begin{remark}
\label{phase:approximate-tangent}
The columns of \(T\) are only numerical directions.  Neither exact
horizontality nor exact orthogonality is assumed.  In particular
\[
 \|T\|\leq\sqrt{1+\mathtt{ort}},
\]
and the explicit \(\mathtt{hor}\) error will be retained when raw
coordinate displacements are projected.  The normal estimates bound
the omitted directions; those directions are not set equal to zero.
No constant-rank or manifold hypothesis on the exact solution set is
needed.
\end{remark}

\subsection{Constructing Root certificates}

The construction uses an auxiliary real map
\(M:\R^{30}\longrightarrow\R^{36}\) to turn the residual equation into
a bound on \((I-TT^*)\delta\).  This map is used only to establish the
Root assertions and is not serialized.  We first control the nonlinear
remainder after applying \(M\).

\begin{lemma}[Quadratic and difference remainders]
\label{phase:remainder}
There is a real \(30\)-vector \(\mathcal R_P(x)\) such that
\[
 F(P+x)=F(P)+L_Px+\mathcal R_P(x).
\]
Put \(a_{ij}(x)=\|x_i-x_j\|^2\), where \(x_i\) is row \(i\) of \(x\).
Then
\begin{align}
 |(\mathcal R_P(x))_{ij}|&\leq\pi a_{ij}(x), \label{phase:edge-remainder}\\
 |(\mathcal R_P(x)-\mathcal R_P(v))_{ij}|
 &\leq\pi\sqrt{a_{ij}(x-v)}
       \bigl(\sqrt{a_{ij}(x)}+\sqrt{a_{ij}(v)}\bigr).
 \label{phase:edge-difference}
\end{align}
Let \(M:\R^{30}\longrightarrow\R^{36}\) be any real matrix.  For real
unit \(y\in\R^{36}\), let \(w_{ij}(y)\) be the Euclidean norm of the two
edge coordinates of \(M^*y\).  If
\begin{equation}
 \sum_{i>j}w_{ij}(y)a_{ij}(x)\leq b\|x\|^2
 \quad\hbox{for every horizontal \(x\) and every unit \(y\),}
 \label{phase:weighted-constant}
\end{equation}
then, for horizontal \(x,v\),
\begin{align}
 \|M\mathcal R_P(x)\|&\leq\pi b\|x\|^2,
 \label{phase:matrix-remainder}\\
 \|M(\mathcal R_P(x)-\mathcal R_P(v))\|
 &\leq\pi b\,\|x-v\|(\|x\|+\|v\|).
 \label{phase:matrix-difference}
\end{align}
\end{lemma}

\begin{proof}
For real \(t\), put
\[
 r(t)=\frac{e(t)-1-2\pi i t}{2\pi i}.
\]
Its derivative is \(e(t)-1\), of modulus at most \(2\pi|t|\).
Integration from zero gives \(|r(t)|\leq\pi t^2\).  On the segment from
\(v\) to \(x\), convexity of absolute value gives
\[
 |r(x)-r(v)|
 \leq2\pi|x-v|\int_0^1|(1-u)v+ux|\,du
 \leq\pi|x-v|(|x|+|v|).
\]
Apply these formulas to each scalar row difference and sum over \(g\).
The unit factors \(e(P_{ig}-P_{jg})\) do not affect modulus.
Cauchy--Schwarz in \(g\) gives \eqref{phase:edge-difference}.

Pairing \(M\mathcal R_P(x)\) with a real unit \(y\) and using the
two-coordinate Cauchy--Schwarz inequality on each edge proves
\eqref{phase:matrix-remainder}.  For the difference use, separately,
\[
 \sum_{i>j}w_{ij}
     \sqrt{a_{ij}(x-v)a_{ij}(x)}
 \leq
 \left(\sum_{i>j}w_{ij}a_{ij}(x-v)\right)^{1/2}
 \left(\sum_{i>j}w_{ij}a_{ij}(x)\right)^{1/2}
 \leq b\|x-v\|\|x\|,
\]
and the same inequality with \(v\) in the last factor.  Take the
supremum over \(y\).  This also proves that the same constant controls
both remainders.
\end{proof}

\begin{lemma}[Two computable remainder constants]
\label{phase:remainder-constants}
Two sufficient choices in \eqref{phase:weighted-constant} are
\[
 b=\sqrt8\,\|M\|,\qquad b=\sqrt6\,\|\mathcal S\|,
\]
where the \(30\)-by-\(216\) real matrix \(\mathcal S\) has row
\[
 \mathcal S_{a,:}
   =M_{:,a}^*\otimes(\mathbf e_i-\mathbf e_j)^*
\]
when \(a\) is either of the two real-coordinate indices of edge
\(i>j\).  Consequently
\(\pi\) times the smaller of any certified upper bounds for these two
quantities is a valid \(\mathtt C\).
\end{lemma}

\begin{proof}
For horizontal \(x\), its row vectors satisfy \(\sum_i x_i=0\).
Let \(s=\|x\|\) and \(u_i=\|x_i\|^2\).  Expansion gives
\begin{align}
 \sum_{i>j}a_{ij}(x)&=6s^2,\nonumber\\
 \sum_{i>j}a_{ij}(x)^2
   &=6\sum_i u_i^2+s^4+2\sum_{i,j}(x_i\cdot x_j)^2.
 \label{phase:row-distance-identities}
\end{align}
Indeed the cross terms use
\(\sum_{j\ne i}x_i\cdot x_j=-u_i\).  Also
\[
 u_i=\left\|-\sum_{j\ne i}x_j\right\|^2
 \leq5(s^2-u_i),\qquad
 \sum_{i,j}(x_i\cdot x_j)^2\leq\sum_{i,j}u_i u_j=s^4.
\]
Thus \(\sum_i u_i^2\leq(5/6)s^4\), and the second sum in
\eqref{phase:row-distance-identities} is at most \(8s^4\).
Cauchy--Schwarz on the edges, with
\(\sum w_{ij}(y)^2=\|M^*y\|^2\), proves the first choice of \(b\).

For the second, apply \(\mathcal S\) to
\(y\otimes x_{:,g}\), then sum squared norms over \(g\).  For unit \(y\)
this gives
\[
 \sum_{i>j}w_{ij}(y)^2a_{ij}(x)
 \leq\|\mathcal S\|^2s^2.
\]
Another Cauchy--Schwarz inequality, now using
\(\sum a_{ij}=6s^2\), proves the claim.  The row Gram of \(\mathcal S\)
is computable without forming its \(216\) columns:
\begin{equation}
 (\mathcal S\mathcal S^*)_{ab}
  =\bigl((\mathbf e_i-\mathbf e_j)\cdot
                (\mathbf e_{i'}-\mathbf e_{j'})\bigr)
           (M_{:,a}\cdot M_{:,b}).
 \label{phase:tensor-gram}
\end{equation}
Here the edge labels are repeated for their real and imaginary parts.
This is the sign-factor formula in \texttt{buildH::patch}.
\end{proof}

\begin{proposition}[Sufficient conditions for a Root certificate]
\label{phase:root-bounds}
Let \(P,T\) and the nonnegative constants be as in
Definition~\ref{phase:root-data}.  Suppose the basic bounds
\eqref{phase:root-basic} hold and an auxiliary stored real
\(36\)-by-\(30\) matrix \(M\) satisfies
\begin{equation}
 \|M\|\leq\mathtt{mnorm},\qquad
 \mathtt{mnorm}\,\mathtt{froot}\leq\mathtt{n0},\qquad
 \|\Pi-TT^*-ML_P\Pi\|\leq\mathtt{id},
 \label{phase:projector-certificate}
\end{equation}
and \(\mathtt C\) is an upper bound for a remainder constant
\(\pi b\) in Lemma~\ref{phase:remainder}.  Then these data form a
certified Root.
\end{proposition}

\begin{proof}
Let \(\delta\) be horizontal with \(F(P+\delta)=0\), and put
\(\eta=T^*\delta\).  Write \(E=\Pi-TT^*-ML_P\Pi\).  The Taylor
equation at the exact solution gives
\[
 (I-TT^*)\delta
  =E\delta-MF(P)-M\mathcal R_P(\delta).
\]
The certificate and Lemma~\ref{phase:remainder} prove
\eqref{phase:root-normal}.  The squared inequality follows without
assuming that \(TT^*\) is a projector:
\[
 \|\delta-T\eta\|^2
   =\|\delta\|^2-2\|\eta\|^2+\eta^*(T^*T)\eta
   \geq\|\delta\|^2-(1+\mathtt{ort})\|\eta\|^2.
\]
In fact this calculation holds for arbitrary \(\delta\).

For any horizontal \(v\), put \(h'=\delta-v\) and subtract the Taylor
formulas at \(\delta\) and \(v\).  Since \(h'\) is horizontal and
\(F(P+\delta)=0\),
\[
 (I-TT^*)h'
   =Eh'-MF(P+v)-M\bigl(\mathcal R_P(\delta)-\mathcal R_P(v)\bigr).
\]
Apply \eqref{phase:matrix-difference}.  This proves
\eqref{phase:root-difference}, including its residual term at the
possibly inexact comparison point \(P+v\).
\end{proof}

The use of a numerical factorization followed by a checked Gram defect
is a standard device in verified matrix-norm estimation; compare
\cite{Rump2011VerifiedSingularValues}. The bounds here are derived from
the displayed residual identities and the source-specific arithmetic
estimates, including the error in the proposed tangent frame.

\begin{proposition}[Certificates produced by the source]
\label{phase:stored-root-certificate}
A successful call to \texttt{buildH::patch}, under the arithmetic
conditions of Section~\ref{sec:rounding}, certifies
\eqref{phase:root-basic}, \eqref{phase:projector-certificate}, and a valid
remainder constant \(\mathtt C\).  Therefore it produces certified Root
data in the sense of Definition~\ref{phase:root-data}.
\end{proposition}

\begin{proof}
The direct checks, rather than the accuracy of the proposed
eigendecomposition, establish the assertion.  The source first proposes
an approximate horizontal basis and a splitting of the \(25\)-coordinate
horizontal Jacobian.  Eigenvalues below \(.64\) propose columns of \(T\);
the other columns propose factors \(H_V,L_V\) for \(M\).
Every accepted factor entry is checked to have magnitude at most \(2\),
and at most \(25\) columns occur.  More than \(10\) proposed \(T\)
columns causes failure.

Certified absolute row sums of the column Grams of \(H_V,L_V\) bound
their squared operator norms.  The bound \(\mathtt{mnorm}\) also
includes the difference between the accumulated stored \(M\) and the
exact product \(H_VL_V^*\) of those stored factors.  The required
additions and the \(2.1\cdot10^{-11}\) Frobenius bound on this difference
are proved in Section~\ref{round:root}.  Thus \(\mathtt{mnorm}\) bounds
the actual stored map, and the first choice in
Lemma~\ref{phase:remainder-constants} supplies a certified \(\mathtt C\).

For the optional tensor improvement the source forms the row Gram
\eqref{phase:tensor-gram}.  It proposes a real change of basis \(U\) and
checks an upper \(\beta\) for \(\|U^*U-I\|\).  If \(\beta<.1\), then
\(\|U^{-1}\|^2\leq1/(1-\beta)\).  A certified absolute-row-sum bound
\(r\) for \(U^*(\mathcal S\mathcal S^*)U\) implies
\[
 \|\mathcal S\|^2
 \leq\frac{r}{1-\beta}.
\]
This uses congruence by an invertible matrix, not assumed exact
orthogonal diagonalization.  The second choice in
Lemma~\ref{phase:remainder-constants}, with the one-sided margins in
Section~\ref{round:root}, gives the alternative \(\mathtt C\).
Both candidates are upper bounds, so their minimum is allowed.

Finally, the source evaluates the full \(36\)-coordinate identity in
\eqref{phase:projector-certificate}, not merely the identity in the
proposed horizontal coordinates.  The padded Frobenius residual gives
\(\mathtt{id}\).  Corresponding checks on \(T^*T-I\), \((I-\Pi)T\),
and \(F(P)\) give the basic bounds, and the stored \(\mathtt{n0}\)
is an upper bound for \(\mathtt{mnorm}\,\mathtt{froot}\).
Section~\ref{round:root} proves the one-sided error estimates for all
these checks and records their acceptance thresholds.  Replacing
operator norms by Frobenius norms only enlarges the bounds.
Proposition~\ref{phase:root-bounds} therefore applies to every successful
call.
\end{proof}

\subsection{Transport and alignment of phase balls}

\begin{lemma}[Equivalence transport]
\label{phase:transport}
Let \(\Phi\) permute rows and columns of a phase matrix, with optional
transpose and optional simultaneous change of all signs.  Then \(\Phi\)
is a real linear isometry commuting with \(\Pi\), and
\begin{equation}
 \|F(\Phi P)\|=\|F(P)\|.
 \label{phase:residual-invariance}
\end{equation}
The residual norm is also invariant under row/column rephasings and
entrywise integer additions.  Replacing \(P,T\) by \(\Phi P,\Phi T\),
with \(\Phi\) applied to every column of \(T\), preserves all the
certified Root assertions with the same scalar constants.  The
existence of an extension by two further mutually unbiased bases is
preserved in both directions.
\end{lemma}

\begin{proof}
Commutation with double centering and preservation of Euclidean norms
follow from the formulas for these maps.  Row or column permutations,
conjugation, and rephasing preserve the residual norm directly.  For
transpose, put \(A=e(P)\).  Both \(AA^*\) and \(A^*A\) have diagonal
entries \(6\), and
\[
 \|AA^*\|_F^2=\|A^*A\|_F^2,\qquad
 \|F(P)\|^2=\frac{\|AA^*\|_F^2-216}{8\pi^2}.
\]
The first identity follows by cyclically permuting the factors under
the trace.  For \(A^T\) the same expression uses \(A^*A\), up to
transpose/conjugation.  This proves
\eqref{phase:residual-invariance} at every real center, not just at a
solution.

Pull back horizontal \(\delta,v\) by \(\Phi^{-1}\) in
\eqref{phase:root-normal}--\eqref{phase:root-difference}.  Inner products,
norms, horizontality, and exact solution status are preserved, and the
residual term is preserved by \eqref{phase:residual-invariance}.
The same argument preserves \eqref{phase:root-basic}.  No transported
matrix \(M\) is required.

Extendibility is preserved by Lemma~\ref{seed:equivalences}.
Its transpose argument applies to each exact solution; it does not
use unitarity of the approximate center.
\end{proof}

Write
\[
 D_0(P)_{ij}=P_{ij}-P_{i0}-P_{0j}+P_{00}.
\]
For any two real centers \(P,Q\) and any integer matrix \(N\),
\begin{equation}
 d=\Pi(D_0(P)-D_0(Q)-N)
 \quad\Longrightarrow\quad
 e(Q+d)\text{ is phase-equivalent to }e(P).
 \label{phase:alignment}
\end{equation}
Indeed \(\Pi D_0(P)=\Pi P\); subtracting \(\Pi\) from the identity
operator leaves a matrix constant on rows plus a matrix constant on
columns.  Thus the difference between \(Q+d\) and \(P-N\) is such a
rephasing.  For every horizontal \(\delta\), the same fixed
alignment makes \(e(Q+d+\delta)\) equivalent to \(e(P+\delta)\).
This is why a particular choice of integers suffices: no optimality
of the wraps is needed.  The source's rounded \texttt{dephase} and
\texttt{wrap} operations determine these integer choices, and
Section~\ref{sec:rounding} bounds their remaining arithmetic error.

\subsection{Phase displacements for a retained completion cell}

Use the notation of Section~\ref{sec:completion}.  In particular \(D\)
is the ideal comparison block at the fixed seed and unitary-cell
center, \(D'\) is the lower-right block of a feasible exact flat
matrix, and \(Z=D'-D\).  The certified bounds and linearization are
\[
 \|Z\|_F\leq z,\qquad
 \|Z-\mathcal J(s,\theta)\|_F\leq E_r,\qquad
 \|s\|\leq R_s,\quad |\theta_a|\leq h\quad(1\leq a\leq4),
\]
as in \eqref{completion:z-bound},
\eqref{completion:linear-map}, and
\eqref{completion:linear-remainder-bound}.  Here \(s\in\R^9\) and
\(\theta\in\R^4\) are in radians.  Denote by \(\mathcal I_s\) the nine
seed positions, whose row-major indices are
\[
 (7,8,9,10,11,13,19,25,31),
\]
and by \(\mathcal I_D\) the \(16\) lower-right positions
\((i,j)\) with \(2\leq i,j\leq5\).  The first row and column have zero
raw displacement.

\begin{lemma}[Arc and weighted-arc bounds]
\label{phase:arc-bounds}
Suppose \(m_j=|D_j|>0\) at every lower-right position and put
\(D_j^\circ=D_j/m_j\).  Choose the shortest angular displacement, in
radians, \(d_j\in[-\pi,\pi]\) from \(D_j^\circ\) to \(D'_j\).  Let
\(\mathtt{slack}\) be a nonnegative upper bound for
\(z^2-\sum_j(m_j-1)^2\), let \(\mathtt{mi}\) be a positive lower bound
for \(\min_jm_j\), and choose
\[
 q_d\geq\frac{\mathtt{slack}}{\mathtt{mi}},\qquad q_d<3.8.
\]
Then valid upper bounds are
\begin{equation}
 \|d\|\leq\mathtt{arc},\qquad
 \left(\sum_jm_jd_j^2\right)^{1/2}\leq\mathtt{weighted},
 \label{phase:arc-weighted}
\end{equation}
provided
\[
 \mathtt{arc}\geq\frac{\sqrt{q_d}}{\sqrt{1-q_d/4}},\qquad
 \mathtt{weighted}\geq
             \frac{\sqrt{\mathtt{slack}}}{\sqrt{1-q_d/4}}.
\]
\end{lemma}

\begin{proof}
Since \(D'_j\) has modulus one,
\begin{equation}
 |D'_j-D_j|^2=(m_j-1)^2+
                         m_j|D'_j-D_j^\circ|^2.
 \label{phase:chord-identity}
\end{equation}
Summing proves
\(\sum_jm_j|D'_j-D_j^\circ|^2\leq\mathtt{slack}\), and hence
\(\sum_j|D'_j-D_j^\circ|^2\leq q_d\).  In particular no individual
chord reaches \(2\).  If \(c_j=|D'_j-D_j^\circ|\), then
\(|d_j|=2\arcsin(c_j/2)\).  On \(0\leq c\leq\sqrt{q_d}\) its
derivative is at most \((1-q_d/4)^{-1/2}\).  Integrating from zero
gives
\(|d_j|\leq c_j(1-q_d/4)^{-1/2}\), and the two summed bounds follow.
\end{proof}

The source's \texttt{Candidate::prep} only attempts this argument when
its direction and magnitude checks pass.  It uses the stored flat
sum \(\mathtt{sum}\) and sets
\[
 \mathtt{slack}
   =\max(0,\mathtt z^2-\mathtt{sum}+10^{-5})+10^{-10},
 \qquad
 q_d=\mathtt{slack}/\mathtt{mi}+10^{-8}.
\]
The subtraction \(10^{-6}\) from each computed magnitude gives
\(\mathtt{mi}\) and the upper reciprocal
\(\mathtt{invmag}_j=1/(\mathtt m_j-10^{-6})+10^{-9}\).
The ratio has \(10^{-7}\) added and the final arc quantities each have
\(10^{-8}\) added.  As proved in Section~\ref{sec:rounding}, these
numbers have the one-sided properties of Lemma~\ref{phase:arc-bounds}
for the ideal comparison block.  Failure of any check means failure
to cover this cell by this route, not an exclusion of the cell.

\begin{lemma}[Stored center and support representation]
\label{phase:center-correction}
Let \(P_0\) be the stored \texttt{Candidate} phase matrix.  On a
successful preparation path there is a correction \(\epsilon\), in
cycles, with
\begin{equation}
 \|\epsilon\|\leq2\cdot10^{-6},
 \label{phase:epsilon}
\end{equation}
such that, up to entrywise integer choices, the actual feasible phase
matrix has the form
\begin{equation}
 P_0+\epsilon+\frac{u}{2\pi}.
 \label{phase:candidate-lift}
\end{equation}
Here the raw angular matrix \(u\) is \(s\) at \(\mathcal I_s\), is
\(d\) at \(\mathcal I_D\), and is zero elsewhere.  In particular
\begin{equation}
 \|u\|\leq\sqrt{R_s^2+\mathtt{arc}^2}.
 \label{phase:raw-displacement}
\end{equation}
Let \(X:\R^{13}\longrightarrow\R^{36}\) give \(s\) at the seed
positions and, at the lower-right positions, give
\begin{equation}
 (X(s,\theta))_j
       =\Im\bigl(\overline{D_j^\circ}\mathcal J(s,\theta)_j\bigr).
 \label{phase:X}
\end{equation}
Then
\begin{equation}
 u=X(s,\theta)+e_D+e_{\sin},\qquad
 \|e_D\|\leq E_r,\qquad
 \|e_{\sin}\|_1\leq\mathtt{arc}^3/6,
 \label{phase:raw-support}
\end{equation}
where both errors are supported on \(\mathcal I_D\).
\end{lemma}

\begin{proof}
The first row and column of \(P_0\) are zero.  Its seed entries are the
binary seed-center phases after
integer reduction; this is the fixed center used in the completion
bounds.  On \(\mathcal I_D\), the source proposes a phase with
\texttt{atan2} but checks its unit direction:
\[
 \bigl|e((P_0)_j)-D^{\rm num}_j/m^{\rm num}_j\bigr|<10^{-7}.
\]
Section~\ref{sec:rounding} bounds the latter normalized direction's
error from the ideal \(D_j^\circ\) by \(5.5\cdot10^{-8}\), and the
total true chord error by \(1.56\cdot10^{-7}\).
For circular phase distance \(t\in[0,1/2]\), the chord
\(2\sin(\pi t)\) is at least \(4t\), by concavity of sine on this
interval.  Thus a suitable integer lift has phase error at most one
quarter of that chord.  Over \(16\) positions this is below
\eqref{phase:epsilon}.  No accuracy property of the proposed
\texttt{atan2} value was assumed.  Adding the shortest block
displacements and the original seed displacements proves
\eqref{phase:candidate-lift}; their supports are disjoint, proving
\eqref{phase:raw-displacement}.

On a lower-right position,
\[
 \Im(\overline{D_j^\circ}Z_j)=\sin d_j.
\]
The difference between this quantity and
\eqref{phase:X} has Euclidean norm at most
\(\|Z-\mathcal J(s,\theta)\|_F\leq E_r\).
Finally
\[
 \sum_j|d_j-\sin d_j|
 \leq\frac16\sum_j|d_j|^3
 \leq\frac16\|d\|^3
 \leq\mathtt{arc}^3/6.
\]
These are the two errors in \eqref{phase:raw-support}.
\end{proof}

\subsection{Direct and tangent bounds for a cell}

Fix a certified Root \(P,T\).  Let \(N\) incorporate the particular
integer choices made by \texttt{Root::finish} and
\texttt{Candidate::delta}, and define the ideal horizontal alignment
\begin{equation}
 \Delta=\Pi(P_0-D_0(P)-N).
 \label{phase:Delta}
\end{equation}
By \eqref{phase:alignment}, every exact feasible matrix in the cell
has an equivalent representation \(e(P+\delta)\), with horizontal
\begin{equation}
 \delta=\Delta+\Pi\left(\epsilon+\frac{u}{2\pi}\right).
 \label{phase:delta}
\end{equation}
The equivalence preserves the extension problem.  The computed
\texttt{delta} array approximates the particular ideal \(\Delta\);
its rounding error is included in the estimates below.

\begin{lemma}[A direct displacement bound]
\label{phase:direct-bound}
Either of the following is an upper bound for \(\Delta\cdot u\):
\begin{align}
 S_1={}&R_s\|(X^*\Delta)_s\|
          +h\|(X^*\Delta)_\theta\|_1
          +E_r\|\Delta_D\|
          +\|\Delta_D\|_\infty\mathtt{arc}^3/6, \label{phase:support-linear}\\
 S_2={}&R_s\|\Delta_s\|
          +\mathtt{weighted}
                  \left(\sum_{j\in\mathcal I_D}\Delta_j^2/m_j\right)^{1/2}.
 \label{phase:support-weighted}
\end{align}
Subscripts \(s,D\) on a phase matrix mean restriction to
\(\mathcal I_s,\mathcal I_D\); subscripts on \(X^*\Delta\) mean its
nine seed and four angle coordinates.  With \(S_\Delta=\min(S_1,S_2)\),
a valid bound is
\begin{equation}
 \|\delta\|\leq
 \left(\|\Delta\|^2+
       \frac{R_s^2+\mathtt{arc}^2}{(2\pi)^2}
       +\frac{S_\Delta}{\pi}\right)^{1/2}
       +\|\epsilon\|.
 \label{phase:s0}
\end{equation}
\end{lemma}

\begin{proof}
Equation~\eqref{phase:raw-support}, the seed ball, and the angle box
give \eqref{phase:support-linear}.  Direct weighted Cauchy--Schwarz
using \eqref{phase:arc-weighted} gives
\eqref{phase:support-weighted}.  Since \(\Delta\) is horizontal,
\(\Delta\cdot\Pi u=\Delta\cdot u\).  Expand the squared norm of
\(\Delta+\Pi u/(2\pi)\), use contraction by \(\Pi\), and then add
\(\|\Pi\epsilon\|\leq\|\epsilon\|\) by the triangle inequality.
\end{proof}

The source's \texttt{covercompute::value} computes these supports and
uses their minimum, with \(+.0001\) before division by \(\pi\).
Its \(\mathtt{d2}\) bounds
\((R_s^2+\mathtt{arc}^2)/(2\pi)^2\), with \(10^{-8}\) added, and
\[
 \mathtt{s0}
   =\sqrt{\mathtt{nd}+\mathtt{d2}
                       +(\mathtt{supp}+.0001)/\pi}+.00001.
\]
Section~\ref{sec:rounding} proves this is an upper bound in
\eqref{phase:s0}: it includes the projection error in
\(\mathtt{nd}\), support-coefficient error, and \(\epsilon\).
An early successful return based solely on this bound is valid.

For a sharper bound put \(\eta=T^*\delta\).  The expression obtained
by temporarily omitting projection and center correction, in radians,
is
\[
 2\pi T^*\Delta+T^*u.
\]
The total omission has the explicit bound
\begin{equation}
 \left\|2\pi\eta-(2\pi T^*\Delta+T^*u)\right\|
 \leq
 \mathtt{hor}\sqrt{R_s^2+\mathtt{arc}^2}
       +2\pi\sqrt{1+\mathtt{ort}}\,\|\epsilon\|.
 \label{phase:tangent-projection-error}
\end{equation}
Indeed \(T^*(I-\Pi)=((I-\Pi)T)^*\), and
\(\|T^*\Pi\epsilon\|\leq\|T\|\|\epsilon\|\).

We next enclose the sum of two linear images of balls by one ellipsoid.
This is a standard outer-sum construction in ellipsoidal calculus; see
\cite[Section~2.2.2, equation~(2.16)]{KurzhanskiyVaraiya2006EllipsoidalToolbox}.
The support-function proof below also covers singular shape matrices.

\begin{lemma}[Ellipsoid containment]
\label{phase:ellipsoid-containment}
For real linear maps \(A,B\), vectors \(\|s\|\leq r\),
\(\|v\|\leq q\), and \(\beta>0\), the sum \(As+Bv\) belongs to
the ellipsoid
\[
 \{G^{1/2}w:\|w\|\leq1\}
 \quad\hbox{whenever}\quad
 G\succeq(1+\beta)r^2AA^*+(1+1/\beta)q^2BB^*.
\]
\end{lemma}

\begin{proof}
For each real test vector \(z\), the support of the sum is at most
\[
 r\|A^*z\|+q\|B^*z\|
 \leq\bigl((1+\beta)r^2\|A^*z\|^2+
                 (1+1/\beta)q^2\|B^*z\|^2\bigr)^{1/2}
 \leq(z^*Gz)^{1/2}.
\]
The scalar inequality follows by completing a square.
For completeness, diagonalizing the positive semidefinite \(G\)
shows that the vectors satisfying
\(|z\cdot x|\leq(z^*Gz)^{1/2}\) for every \(z\) are exactly
\(G^{1/2}\) times the unit ball: zero-eigenvalue directions force
zero components, and taking
\(z_a=x_a/\lambda_a\) in the positive-eigenvalue coordinates gives
\(\sum x_a^2/\lambda_a\leq1\).
This proves the claimed containment, including singular \(G\).
\end{proof}

Partition \(X=(X_s,X_\theta)\), and let \(T_s,T_D\) denote the rows
of \(T\) restricted to the indicated position sets.  From the angle
box, \(2\pi T^*\Delta+hT^*X_\theta[-1,1]^4\) is the convex hull of
the \(16\) offsets
\begin{equation}
 U_\sigma=2\pi T^*\Delta+
                h\sum_{a=1}^{4}\sigma_a T^*X_{\theta,a},
 \qquad \sigma_a\in\{-1,1\}.
 \label{phase:offset-corners}
\end{equation}
By Lemma~\ref{phase:ellipsoid-containment} with \(\beta=.45\), the
sum of the seed and \(e_D\) contributions to \(T^*u\) lies in an
ellipsoid whenever
\begin{equation}
 G_{\rm lin}\succeq
    1.45R_s^2(T^*X_s)(T^*X_s)^*
       +(1+1/.45)E_r^2T_D^*T_D.
 \label{phase:ellipsoid-linear}
\end{equation}
The remaining sine error has norm at most
\[
 \left(\max_{j\in\mathcal I_D}\|T_{j,:}\|\right)
                               \mathtt{arc}^3/6.
\]
There is also a bound using the raw seed and weighted angle balls,
with a single offset \(U_0=2\pi T^*\Delta\):
\begin{equation}
 G_{\rm dir}\succeq
    1.8R_s^2T_s^*T_s+
       (1+1/.8)\mathtt{weighted}^2
                     T_D^*\operatorname{diag}(1/m_j)T_D.
 \label{phase:ellipsoid-direct}
\end{equation}
This is Lemma~\ref{phase:ellipsoid-containment} with \(\beta=.8\).
It has no sine-error term because it bounds \(d\) directly.

\begin{lemma}[Offset ellipsoid norm bound]
\label{phase:ellipsoid-norm}
Let \(G\succeq0\) and \(U\in\R^k\).  For every
\(\gamma>\lambda_{\max}(G)\) and \(\|w\|\leq1\),
\begin{equation}
 \|U+G^{1/2}w\|^2
   \leq\gamma\bigl(1+U^*(\gamma I-G)^{-1}U\bigr).
 \label{phase:ellipsoid-gamma}
\end{equation}
The simpler upper bound
\(\|U\|+\sqrt{\lambda_{\max}(G)}\) for the norm is also valid.
If offsets range over a convex hull of finitely many \(U_\sigma\),
it suffices to take the largest of the bounds at those offsets.
\end{lemma}

\begin{proof}
Put \(A=\gamma I-G\succ0\).  Expanding and completing the square,
\[
 \|U+G^{1/2}w\|^2
 \leq\gamma+\|U\|^2+
       2w^*G^{1/2}U-w^*Aw
 \leq\gamma+\|U\|^2+
       U^*G^{1/2}A^{-1}G^{1/2}U.
\]
Since \(G\) commutes with \(A^{-1}\), the last two terms equal
\(\gamma U^*A^{-1}U\).  The simpler bound is the triangle
inequality.  For fixed \(w\), convexity of norm bounds the norm at
any convex combination of offsets by the maximum over those
offsets; taking the supremum over \(w\) proves the last assertion.
\end{proof}

Let \(\mathcal E(G;\{U_\sigma\})\) denote any certified upper bound
obtained from Lemma~\ref{phase:ellipsoid-norm}.  The preceding
arguments give either of the tangent bounds
\begin{align}
 \|T^*\delta\|
 \leq{}&
 \frac{\mathcal E(G_{\rm lin};\{U_\sigma\})+
      \max_{j\in\mathcal I_D}\|T_{j,:}\|\,\mathtt{arc}^3/6+
      \mathtt{hor}\sqrt{R_s^2+\mathtt{arc}^2}}{2\pi}
       +\sqrt{1+\mathtt{ort}}\|\epsilon\|,
 \label{phase:tangent-linear}\\
 \|T^*\delta\|
 \leq{}&
 \frac{\mathcal E(G_{\rm dir};\{U_0\})+
      \mathtt{hor}\sqrt{R_s^2+\mathtt{arc}^2}}{2\pi}
       +\sqrt{1+\mathtt{ort}}\|\epsilon\|.
 \label{phase:tangent-direct}
\end{align}

\paragraph{Numerical realization of these tangent bounds.}
In \texttt{covercompute}, \texttt{off} is the computed
\(2\pi T^*\Delta\), \texttt{AB} is the computed \(T^*X\), and
\texttt{U} holds the \(16\) corners.  The covariance calculations use
the \emph{stored} \texttt{AB} coefficients; their errors from the
ideal coefficients are charged separately outside the ellipsoid.
More explicitly, if \(\widehat A\) denotes the stored
\texttt{AB} matrix, the first stored covariance dominates the
right-hand side of \eqref{phase:ellipsoid-linear} with
\(T^*X_s\) replaced by \(\widehat A_s\), and its corners use
\(\widehat A_\theta\).  Section~\ref{sec:rounding} gives
\[
 \|\widehat A-T^*X\|<1.5\cdot10^{-5},\qquad
 \|(s,\theta)\|<.226,
\]
so this substitution costs less than \(3.4\cdot10^{-6}\) radians.
For both covariance choices a diagonal \(10^{-7}I\) is added.  The
reciprocals in the second choice are the certified upper
\texttt{invmag} values, so this change also enlarges the covariance.
Section~\ref{sec:rounding} proves positive semidefinite domination of
the exact covariance expressions in this stored-coefficient model.

The function \texttt{spectralUpper} bounds the norm of a symmetric
stored \(G\) by the smaller of a maximum absolute row sum plus
\(10^{-10}\), and
\(\|G^2\|_F^{1/2}\) computed with \(2\cdot10^{-7}\) added.  Thus
\texttt{ell} starts with a certified triangle-inequality bound.
For each attempted \(\gamma\), it uses a lower gap
\(\mathtt{gap}\leq\gamma-\lambda_{\max}(G)\), requiring
\(\mathtt{gap}\geq10^{-6}\).  If \(z\) is any finite trial solve and
\(r=U-(\gamma I-G)z\), then
\begin{equation}
 U^*(\gamma I-G)^{-1}U
 \leq U^*z+\frac{\|U\|\|r\|}{\mathtt{gap}}.
 \label{phase:inverse-residual}
\end{equation}
This follows by substituting
\((\gamma I-G)^{-1}U=z+(\gamma I-G)^{-1}r\).
The source checks the finite scale of \(z\), bounds this residual
expression with its explicit \(10^{-10}\)-scale additions, and adds
\(10^{-6}\) after taking the candidate square root and again upon
return.  It does not rely on accuracy of its Cholesky solve.
Section~\ref{sec:rounding} verifies the errors including division by
the lower gap.

On a potentially successful completion-cell path, the external
error in either applicable tangent bound, in radians, is below
\(.00012\).  This includes the coefficient and offset errors, the
projection term in \eqref{phase:tangent-projection-error} (less than
\(8.81\cdot10^{-5}\) radians), and the center correction (less than
\(1.27\cdot10^{-5}\) radians).
Accordingly the source adds \(.00015\) in radians before division
by \(2\pi\), and another \(.00002\) in cycles afterward.  The sine
term is included only for \eqref{phase:tangent-linear}.  These
margins establish an upper \(t_0\geq\|T^*\delta\|\) while retaining
both the phase lift and approximate-horizontality errors.

\subsection{Bootstrap and local coverage}

\begin{proposition}[Finite improvement of a certified radius]
\label{phase:bootstrap}
Suppose \(s_0\geq\|\delta\|\) and
\(t_0\geq\|T^*\delta\|\) for every feasible solution in a completion
cell, using the particular alignment \eqref{phase:delta}.  Starting
with any such upper \(s\), a further upper is
\begin{equation}
 \left(n_P(s)^2+(1+\mathtt{ort})t_0^2\right)^{1/2}.
 \label{phase:bootstrap-update}
\end{equation}
Replacing \(s\) by this number only when it decreases \(s\), and
stopping after any finite number of steps, preserves coverage of
every feasible solution in the cell.
\end{proposition}

\begin{proof}
The constants in \(n_P\) are nonnegative, so
\(n_P(\|\delta\|)\leq n_P(s)\).  Substitute
\eqref{phase:root-normal} into \eqref{phase:root-square} and use
the given tangent upper.  Each proposed update is an upper bound
at the current stored upper, so induction proves the assertion.
No convergence or uniqueness assertion is involved.
\end{proof}

The source's \texttt{bootstrap} uses at most \(30\) updates.  It
inflates its tangent input by \(.00002\) and its initial radius by
\(10^{-6}\), then evaluates \eqref{phase:bootstrap-update} with
\(10^{-8}\) added at each attempted update.  The positive-expression
error bounds of Section~\ref{sec:rounding} make each update an upper.
When \(k=0\), the mathematical tangent norm is zero; the source's
small positive tangent input is harmless.

Scale failures in \texttt{ell} return \(100\), which makes that
tangent path fail \texttt{bootstrap}'s \(t<10\) guard.  Its failure
sentinel is \(10\).  These values cannot pass the local target
\(.103\); they are not certificates.  The other guard requires
\(0\leq\mathtt{init}<3\) and \(\mathtt C<20\).  A direct valid
\(\mathtt{s0}\) can still establish coverage without either
ellipsoid path.

\texttt{local::test} returns success only after a certified
\texttt{covercompute} bound for a specific Root fits the target with
its final margins.  It grows that Root's recorded radius to cover
the bound, adding \(10^{-8}\).  Selecting only a few nearby Roots
by a heuristic score cannot invalidate this implication.  If
those fail, snapping merely proposes a new center, which must pass
\texttt{buildH::patch} and the same covering computation.  Failed
snapping or failed coverage does not discard the completion cell:
the completion search subdivides it, or records an unresolved
seed if its depth or node limit is reached, as described in
Section~\ref{sec:completion}.

\subsection{Merging the global atlas}

\begin{definition}[Record and its covered solution set]
\label{phase:record}
A \texttt{Record} is a certified Root with a nonnegative radius \(r\).
Its solution set consists of matrices \(e(P+\delta)\) with
\(\delta\) horizontal, \(\|\delta\|\leq r\), and \(F(P+\delta)=0\).
Coverage by Records may be up to the extendibility-preserving
equivalences of Lemma~\ref{phase:transport}; the chosen equivalence
may differ for different assignments.  The term \emph{sphere} in the
source includes the interior and boundary of this closed ball.
\end{definition}

\begin{lemma}[A certified atlas merge]
\label{phase:atlas-merge}
Let a local Record have center \(P\) and radius \(r\).  Let \(Q\)
be an existing certified atlas center.  Suppose an allowed
\(\Phi\) and integer matrix \(N\) have been selected, and \(d\) is
a certified upper for
\[
 \left\|\Pi\bigl(D_0(\Phi P)-D_0(Q)-N\bigr)\right\|.
\]
Then the local Record is covered up to equivalence by the Record
at \(Q\) of radius \(r+d\).
\end{lemma}

\begin{proof}
An exact solution \(e(P+\delta)\) from the local Record first
transports to \(e(\Phi P+\Phi\delta)\).  The center alignment
\eqref{phase:alignment} represents it at \(Q\) by the horizontal
displacement
\[
 \Pi(D_0(\Phi P)-D_0(Q)-N)+\Phi\delta.
\]
Its norm is at most \(d+r\).  Both transformations preserve exact
solution status and extendibility.
\end{proof}

The implementation in \texttt{atlas.cpp} uses \texttt{mapH} for
\(\Phi\).  Its two index arrays are genuine permutations.  To see
this directly in \texttt{variants}, the leading two row indices
are distinct and \texttt{complement} supplies precisely the other
four.  The six column labels remain attached to six distinct
integers when their phase keys are sorted and cyclically reordered.
The later tail sort and rotation only permute the four remaining
row labels.  With the transpose flag the lookup is
\(6\,\mathtt{col}[j]+\mathtt{row}[i]\), rather than
\(6\,\mathtt{row}[i]+\mathtt{col}[j]\); either lookup is a bijection
of all \(36\) entries.  The sign flag multiplies every entry, and
every column of \(T\), by the same sign.

For a finite certified center, \texttt{canonicalMap} obtains at
least one callback, so it does not use an uninitialized or
degenerate index choice.  In each finite pivot list a largest gap
is attained and is admitted by the nonnegative tolerance.  A list
attaining the smallest such largest gap is admitted by the next
comparison.  Among its five tail scores a maximum is attained
and is admitted by the final comparison.  The first resulting
callback initializes the chosen map.  Ties and repeated phase
values do not alter the distinct attached labels.  The signature
used to choose among these maps need not be a canonical invariant.

\texttt{Atlas::find} computes an explicit wrapped dephased
difference, projects it, and evaluates its norm.  On a successful
return it adds \(10^{-7}\) to the computed norm, exceeding the
dephasing/projection errors certified in
Section~\ref{sec:rounding}.  \texttt{Atlas::add} then grows the
selected radius to at least the incoming radius plus this
distance, with another \(10^{-8}\) added and a separate
\(2\cdot10^{-8}\)-margin comparison to the cap.  This is
Lemma~\ref{phase:atlas-merge}.  It remains valid when an alternative
map found by \texttt{variants} supplied the distance; only existence
of that allowed map is needed to establish the merge.

If no merge is certified, the source inserts the transformed
Root and its original local radius.  The hash buckets, pivot
selection, and early raw-distance cuts only reduce the
opportunities to merge; none removes an input Record.  The A
atlas uses radius cap \(.125\).  Its size cap aborts rather than
certifying a discarded Record.

\begin{proposition}[Output meaning of a successful A stage]
\label{phase:A-cover}
Suppose the initial seed cover is valid and the completion search
processes all its seeds with no failed or unresolved seed.  Then
the A-stage Records cover, up to extendibility-preserving
equivalence, every exact second-basis Hadamard matrix required
for the four-basis problem.
\end{proposition}

\begin{proof}
The completion parametrization covers each feasible matrix in a
seed by its closed unitary cells.  Follow any such matrix through
containing cells.  A sound exclusion cannot reject it; subdivision
retains a containing child.  Since the search is finite and the
failed-setup and unresolved depth/node-limit paths are assumed
absent, this path must end with a certified local covering bound.
The preceding propositions give a certified Root ball for that
assignment.  Atlas insertions preserve these
balls, and every merge preserves their solution sets up to
equivalence by Lemma~\ref{phase:atlas-merge}.  This proves the
claim.  Boundaries cause no exception: all seed balls, unitary
cells, ellipsoids, and phase balls used in this argument may be
taken closed, and overlapping assignments are harmless.
\end{proof}
